\chapter{Contexte général du projet}
\label{chap:contexte}

\section{Présentation générale}
Ce chapitre établit les fondations du projet \devpilot{}. Il caractérise le
besoin métier, analyse l'existant à travers une étude comparative des solutions
concurrentes, puis formalise les exigences fonctionnelles et non fonctionnelles,
les contraintes, le cahier des charges et la gestion des risques. Cette analyse
conditionne l'ensemble des décisions de conception et d'architecture présentées
dans les chapitres~\ref{chap:conception} et~\ref{chap:archi}.

\subsection{Contexte}
Les équipes d'ingénierie logicielle produisent et consomment en permanence une
connaissance technique abondante : spécifications, documents d'architecture, notes
de conception, comptes rendus, directives internes et documentation de projet.
Cette connaissance possède trois propriétés qui la rendent difficile à exploiter.
Elle est d'abord \textbf{volumineuse et croissante}. Elle est ensuite
\textbf{fragmentée}, répartie entre des outils, des formats et des silos
hétérogènes. Elle est enfin \textbf{contextuelle} : une même question n'appelle
pas la même réponse selon le projet, l'équipe ou la version concernée.

La recherche par mots-clés des outils traditionnels échoue dès que le vocabulaire
de la question diffère de celui du document ; elle ne comprend pas l'intention
sémantique et renvoie des listes de documents que l'utilisateur doit encore lire
et synthétiser. À l'inverse, les assistants génératifs généralistes produisent
une synthèse directe, mais répondent à partir de paramètres figés au moment de
l'entraînement, sans connaissance du patrimoine documentaire propre à
l'organisation et sans garantie de véracité.

\subsection{Organisme d'accueil}
Ce projet de fin d'études a été réalisé au sein de [\textit{nom de l'organisme
d'accueil — à compléter}], [\textit{secteur d'activité — à compléter}]. La
structure d'accueil [\textit{présentation succincte de l'organisme, de ses
activités et de l'équipe d'encadrement — à compléter}]. Le projet s'inscrit dans
[\textit{contexte de la mission — à compléter}].

\subsection{Problématique}
La problématique centrale se décompose en deux familles de limites
complémentaires. D'une part, les moteurs de recherche internes, fondés sur
l'appariement lexical, ne comprennent pas la sémantique des requêtes. D'autre
part, les grands modèles de langage généralistes souffrent
d'\textbf{hallucinations}, d'une \textbf{absence de traçabilité}, d'une
\textbf{absence de citations}, d'une \textbf{absence de visibilité} sur la qualité
de récupération et d'une \textbf{absence d'observabilité}~\cite{lewis2020}.

La question directrice du projet est donc la suivante : \emph{comment concevoir et
réaliser une plateforme d'IA capable de centraliser les connaissances d'une équipe
d'ingénierie, d'en récupérer l'information pertinente, de générer des réponses
fondées et citées, d'en mesurer la qualité, et d'offrir une traçabilité et une
observabilité permettant d'auditer et d'exploiter le système comme tout service
critique d'entreprise~?}

\subsection{Objectifs}
Le projet vise à concevoir et implémenter \devpilot{}, une plateforme agentique de
\emph{Génération Augmentée par Récupération} (RAG) multi-espaces. Les objectifs
opérationnels sont synthétisés dans le tableau~\ref{tab:objectifs}.

\begin{table}[H]
\centering\caption{Objectifs opérationnels et capacités correspondantes}
\label{tab:objectifs}\small
\begin{tabularx}{\textwidth}{@{}>{\bfseries}l X@{}}
\toprule
Objectif & Capacité réalisée \\
\midrule
Centraliser & Espaces de travail isolés, gestion documentaire, ingestion asynchrone \\
Récupérer & Recherche sémantique et hybride, reclassement par pertinence \\
Générer & Réponses ancrées produites par un pipeline agentique de sept agents \\
Citer & Génération de citations reliées aux passages sources \\
Mesurer & Évaluation de la fidélité, de la pertinence et du risque d'hallucination \\
Diffuser & \emph{Streaming} temps réel des réponses (Server-Sent Events) \\
Superviser & RAGOps : santé des pipelines, couverture, file de traitement, alertes \\
Tracer & Explorateur et inspecteur de traces de bout en bout \\
Visualiser & AI Execution Studio : visualisation animée du pipeline d'exécution \\
\bottomrule
\end{tabularx}
\end{table}

\subsection{Méthodologie}
La démarche adoptée est itérative et orientée production. Le développement s'est
déroulé en \textbf{vingt-trois phases} incrémentales, de la fondation
conteneurisée jusqu'aux fonctionnalités phares de visualisation. Chaque phase fixe
un objectif vérifiable, introduit des fonctionnalités, justifie ses choix
techniques et est validée par des tests dédiés. Cette progression, détaillée tout
au long du rapport, garantit que le système reflète une évolution réelle et
maîtrisée plutôt qu'une conception monolithique. Le chapitre~\ref{chap:archi}
constitue le c{\oe}ur du travail d'ingénierie.

\subsection{Analyse de l'existant}
Aucune solution unique du marché ne couvre simultanément l'ancrage documentaire
privé, la traçabilité et l'observabilité dans un cadre auto-hébergeable. Nous
analysons ci-dessous six familles de solutions représentatives.

\paragraph{ChatGPT (et assistants génériques).} \textbf{Forces :} qualité
conversationnelle, fluidité, large couverture générale. \textbf{Limites :} répond
depuis ses paramètres, pas depuis la base de connaissances privée ; hallucinations ;
absence de citations vérifiables. \textbf{Manques :} traçabilité, isolation des
données, observabilité opérationnelle.

\paragraph{LangSmith.} \textbf{Forces :} excellente observabilité et traçabilité
des applications à base de LLM~\cite{langsmith}. \textbf{Limites :} c'est un outil
d'\emph{observabilité}, non une plateforme RAG complète et auto-suffisante.
\textbf{Manques :} moteur d'ingestion/récupération intégré, interface utilisateur
métier, gouvernance multi-espaces.

\paragraph{LlamaIndex.} \textbf{Forces :} \emph{framework} riche pour
l'indexation et la récupération~\cite{llamaindex}. \textbf{Limites :} bibliothèque
de développement, non un produit clé en main ; pas d'observabilité ni d'interface
intégrées. \textbf{Manques :} RAGOps, traçabilité persistée, expérience
utilisateur complète.

\paragraph{OpenWebUI.} \textbf{Forces :} interface conversationnelle agréable,
auto-hébergeable. \textbf{Limites :} centrée sur le \emph{chat}, observabilité et
évaluation limitées. \textbf{Manques :} évaluation systématique de la qualité,
traçabilité fine, supervision de la santé du pipeline.

\paragraph{Haystack.} \textbf{Forces :} \emph{framework} RAG modulaire et
mature~\cite{haystack}. \textbf{Limites :} orienté développeur ; l'observabilité
et l'interface restent à construire. \textbf{Manques :} tour de contrôle
opérationnelle prête à l'emploi, visualisation pédagogique de l'exécution.

\paragraph{Plateformes de connaissance d'entreprise.} \textbf{Forces :}
intégration documentaire et gouvernance. \textbf{Limites :} recherche souvent
lexicale ; capacités génératives et de traçabilité du raisonnement limitées.
\textbf{Manques :} pipeline RAG agentique transparent et évalué.

\paragraph{Différenciation de \devpilot{}.} \devpilot{} se distingue en
\textbf{intégrant} l'application RAG \emph{et} son exploitation dans un système
unique, auto-hébergeable : base de connaissances privée par espace de travail,
réponses ancrées et citées, évaluation automatique, traçabilité complète
(RAGOps, explorateur et inspecteur de traces) et visualisation d'exécution
(AI Execution Studio). La transparence du pipeline de raisonnement est le c{\oe}ur
de la proposition de valeur.

\subsection{Étude comparative}
Le tableau~\ref{tab:comparaison} positionne \devpilot{} face aux solutions
analysées selon les capacités déterminantes pour une plateforme d'IA d'entreprise.

\begin{table}[H]
\centering\caption{Étude comparative des solutions}\label{tab:comparaison}
\scriptsize\renewcommand{\arraystretch}{1.25}
\begin{tabularx}{\textwidth}{@{}>{\raggedright\arraybackslash}p{3.1cm} *{6}{>{\centering\arraybackslash}X}@{}}
\toprule
\textbf{Capacité} & \textbf{ChatGPT} & \textbf{LangSmith} & \textbf{LlamaIndex} & \textbf{Haystack} & \textbf{OpenWebUI} & \textbf{\devpilot{}} \\
\midrule
Réponses ancrées sur vos documents & Limité & N/A & Oui & Oui & Partiel & \textbf{Oui} \\
Citations vérifiables & Non & N/A & Partiel & Partiel & Partiel & \textbf{Oui} \\
Traçabilité complète des requêtes & Non & Oui & Non & Non & Non & \textbf{Oui} \\
Isolation par espace de travail & Non & N/A & Non & Non & Partiel & \textbf{Oui} \\
Score d'hallucination & Non & Partiel & Non & Non & Non & \textbf{Oui} \\
Métriques de récupération & Non & Oui & Partiel & Partiel & Non & \textbf{Oui} \\
Chronologie d'exécution des agents & Non & Oui & Non & Non & Non & \textbf{Oui} \\
Auto-hébergement & Non & Partiel & Oui & Oui & Oui & \textbf{Oui} \\
Évaluation intégrée & Non & Partiel & Non & Partiel & Non & \textbf{Oui} \\
Observabilité opérationnelle (RAGOps) & Non & Partiel & Non & Non & Non & \textbf{Oui} \\
\bottomrule
\end{tabularx}
\end{table}

\subsection{Besoins fonctionnels}
Les exigences fonctionnelles (EF) sont synthétisées dans le
tableau~\ref{tab:ef}.

\begin{table}[H]
\centering\caption{Exigences fonctionnelles principales}\label{tab:ef}\small
\begin{tabularx}{\textwidth}{@{}l l X@{}}
\toprule
\textbf{Réf.} & \textbf{Domaine} & \textbf{Exigence} \\
\midrule
EF-1 & Authentification & Inscription, connexion et sessions sécurisées par jetons JWT \\
EF-2 & Autorisation & Contrôle d'accès par rôles (utilisateur, admin, super-admin) \\
EF-3 & Espaces & Création d'espaces isolés et gestion des membres \\
EF-4 & Documents & Téléversement, listing, suivi de statut et suppression \\
EF-5 & Ingestion & Traitement asynchrone : extraction, découpage, vectorisation, indexation \\
EF-6 & Récupération & Recherche sémantique et hybride sur la base vectorielle \\
EF-7 & Reclassement & Reclassement des candidats par pertinence (\emph{reranking}) \\
EF-8 & Génération & Production de réponses ancrées par un pipeline agentique \\
EF-9 & Citations & Génération de citations reliées aux passages utilisés \\
EF-10 & Conversation & Historique de conversations et de messages persistants \\
EF-11 & Streaming & Diffusion incrémentale des réponses (jeton par jeton) \\
EF-12 & Évaluation & Mesure de la fidélité, de la pertinence et du risque d'hallucination \\
EF-13 & Correction & Réparation automatique des réponses non fidèles \\
EF-14 & Administration & Tableaux de bord d'administration et journal d'audit \\
EF-15 & RAGOps & Supervision de la santé des pipelines et de la file de traitement \\
EF-16 & Traçabilité & Explorateur et inspecteur de traces de requêtes RAG \\
EF-17 & Visualisation & AI Execution Studio : visualisation animée de l'exécution \\
\bottomrule
\end{tabularx}
\end{table}

\subsection{Besoins non fonctionnels}
\begin{description}[leftmargin=2.4cm,style=nextline]
\item[Sécurité] Authentification JWT, autorisation par rôles, isolation stricte
des données par espace, journalisation d'audit et masquage systématique des
secrets dans les vues de débogage.
\item[Performance] API asynchrone non bloquante, déport du traitement lourd via
Celery, recherche par plus proches voisins approchés (HNSW) et \emph{streaming}
réduisant le \emph{time-to-first-token}.
\item[Scalabilité] \textit{Backend} sans état horizontalement extensible,
\emph{workers} Celery ajoutables, base vectorielle Qdrant partitionnable.
\item[Disponibilité] Conteneurisation, dépendances de santé explicites et
politiques de relance des tâches asynchrones.
\item[Maintenabilité] Architecture en couches, typage statique de bout en bout,
migrations de schéma versionnées (Alembic).
\item[Observabilité] Traçabilité complète, métriques Prometheus, tableaux de bord
Grafana, RAGOps et inspecteur de traces.
\item[Fiabilité] Repli de modèles de génération, agent correcteur et relance
automatique des ingestions échouées.
\item[Utilisabilité] Interface conversationnelle réactive, rendu Markdown, cartes
de citation interactives, thème clair/sombre.
\end{description}

\subsection{Contraintes}
\begin{description}[leftmargin=2.4cm,style=nextline]
\item[Techniques] Reproductibilité par conteneurs ; possibilité d'exécution
hors-ligne via des modèles locaux (Ollama) ; budget mémoire et calcul maîtrisé.
\item[Académiques] Démonstration d'une démarche d'ingénierie rigoureuse et d'un
raisonnement architectural.
\item[Temps] Réalisation dans le cadre temporel d'un projet de fin d'études.
\item[Ressources] Environnement d'exécution mono-machine, sans infrastructure
\emph{cloud} dédiée.
\end{description}

\subsection{Cahier des charges}
Le tableau~\ref{tab:cdc} relie les objectifs aux livrables vérifiables.

\begin{table}[H]
\centering\caption{Synthèse du cahier des charges}\label{tab:cdc}\small
\begin{tabularx}{\textwidth}{@{}>{\bfseries}l X@{}}
\toprule
Objectif & Livrable \\
\midrule
Ingestion & Pipeline asynchrone Celery (extraction, découpage, vectorisation, indexation Qdrant) \\
Récupération & Service de recherche hybride avec reclassement \\
Génération & Flux agentique de sept agents avec citations \\
Évaluation & Module d'évaluation (fidélité, pertinence, hallucination) \\
Exploitation & Tour de contrôle RAGOps \\
Traçabilité & Explorateur et inspecteur de traces \\
Visualisation & AI Execution Studio \\
Déploiement & Orchestration Docker Compose (8 services) \\
\bottomrule
\end{tabularx}
\end{table}

\subsection{Gestion des risques}
Le tableau~\ref{tab:risques} recense les principaux risques et leurs mesures de
mitigation.

\begin{table}[H]
\centering\caption{Gestion des risques}\label{tab:risques}\small
\begin{tabularx}{\textwidth}{@{}X c c X@{}}
\toprule
\textbf{Risque} & \textbf{Prob.} & \textbf{Impact} & \textbf{Mitigation} \\
\midrule
Hallucination du LLM & Moyenne & Élevé & Ancrage strict, évaluation de fidélité, agent correcteur \\
Indisponibilité du fournisseur LLM & Faible & Élevé & Modèle de repli (\emph{fallback}) ; embeddings locaux \\
Couverture d'\emph{embedding} incomplète & Moyenne & Moyen & Suivi RAGOps de la couverture, relance d'ingestion \\
Échec d'ingestion documentaire & Moyenne & Moyen & Relance Celery (3 tentatives), centre des échecs \\
Volume de \emph{node\_modules} conteneur & Faible & Faible & Procédure de réinstallation documentée \\
Dégradation des performances & Faible & Moyen & Recherche HNSW, \emph{streaming}, chargement différé \\
\bottomrule
\end{tabularx}
\end{table}

\subsection{Conclusion}
Ce chapitre a établi que le besoin n'est pas un \emph{chatbot} supplémentaire,
mais une plateforme d'IA \emph{fondée}, \emph{mesurable} et \emph{observable}.
L'étude comparative montre que \devpilot{} occupe une position distinctive en
unifiant moteur RAG, exploitation et traçabilité. Les exigences non fonctionnelles
— sécurité, scalabilité, observabilité — sont aussi structurantes que les
exigences fonctionnelles. Le chapitre suivant traduit ces besoins en modèles de
conception.
